\documentclass[10pt,a4paper]{article}
\usepackage[T1]{fontenc}
\usepackage{lmodern}
\usepackage[margin=23mm]{geometry}
\usepackage{amsmath,amssymb,amsthm,mathtools}
\usepackage[hidelinks]{hyperref}
\newtheorem{theorem}{Theorem}
\newtheorem{lemma}{Lemma}
\newtheorem{proposition}{Proposition}
\theoremstyle{definition}
\newtheorem{definition}{Definition}
\theoremstyle{remark}
\newtheorem{remark}{Remark}
\newcommand{\R}{\mathbb R}
\newcommand{\Ss}{\mathbb S}
\newcommand{\COP}{\operatorname{COP}}
\newcommand{\col}{\operatorname{col}}
\newcommand{\diag}{\operatorname{diag}}
\newcommand{\conv}{\operatorname{conv}}
\newcommand{\cone}{\operatorname{cone}}
\newcommand{\sat}{\operatorname{sat}}
\newcommand{\mc}{\mathcal M_c}
\newcommand{\Fenv}{\mathcal F_{\underline\Sigma,\overline\Sigma}}
\allowdisplaybreaks
\setlength{\parindent}{0pt}
\setlength{\parskip}{5pt}
\title{Quadratic and Incremental Quadratic Constraints for Saturations\texorpdfstring{\\}{ }and Functions Between Saturation Envelopes}
\author{}
\date{17 November 2026}
\begin{document}
\maketitle

\begin{abstract}
Using the notation of Vahedi Noori et al.~\cite{relu}, this note first derives static QCs, incremental QCs, and the resulting hard IQCs and IIQCs for fixed saturation maps. It then treats the class of functions bounded by two fixed saturation envelopes. Polyhedral graph generators give complete copositive characterizations; a ReLU lift gives implementable sufficient families. A strict gap between the envelopes makes the unrestricted class too broad to admit informative homogeneous incremental QCs. The final construction adds the slope restriction available for sigmoids. These are new derivations in this note, not saturation theorems asserted in~\cite{relu}.
\end{abstract}

\section{Notation and the type of IQC considered}
Write $n=n_v$ and use $\Ss^n$, $D^n$, $DH^n$, and $\COP^n$ as in~\cite{relu}. The scalar ReLU is $\phi(v)=\max\{0,v\}$ and $\Phi_m:\R^m\to\R^m$ is its repeated version. A subscript on $\Phi$ only specifies its dimension. Denote the paper's complete ReLU multiplier cones by $\mc(\Phi_m)$ and $\mc^{inc}(\Phi_m)$.

As in Definitions 1--2 of~\cite{relu}, $F\in QC(M)$ means
\begin{equation}\label{eq:qc}
 \col(v,F(v))^T M\col(v,F(v))\ge0\quad\forall v,
 \qquad M\in\Ss^{2n}.
\end{equation}
Similarly, $F\in QC^{inc}(M)$ means
\begin{equation}\label{eq:incqc}
 \col(\bar v-\hat v,F(\bar v)-F(\hat v))^T M
 \col(\bar v-\hat v,F(\bar v)-F(\hat v))\ge0
 \quad\forall\bar v,\hat v.
\end{equation}
For a function class, membership means that the inequality holds for every fixed function in that class. The same function is used at $\bar v$ and $\hat v$.

Because saturation has nonzero breakpoints and offsets, also allow an explicitly identified \emph{affine lift}
\begin{equation}\label{eq:affqc}
 x=\col(v,w,1),\qquad x^TNx\ge0,\qquad N\in\Ss^{2n+1}.
\end{equation}
Let $E=[I_{2n}\ \ 0]$. The paper's homogeneous QC is the special case $N=E^TME$. A result for $N$ is not silently identified with one for $M$.

For the memoryless operator $w(t)=F(v(t))$, every pointwise QC immediately gives a \emph{hard static IQC}
\begin{equation}\label{eq:hard}
 \int_0^T \col(v(t),w(t))^TM\col(v(t),w(t))\,dt\ge0
 \quad\forall T\ge0.
\end{equation}
Every incremental QC gives a hard static IIQC by replacing $(v,w)$ by $(\bar v-\hat v,\bar w-\hat w)$. In the paper's discrete-time setting, replace the integral by $\sum_{k=0}^{N}$. These factorizations have $\Psi=I$ and $\Pi=M$; no nontrivial dynamic filter is claimed. Affine lifts give finite-horizon inequalities on $\col(v,w,1)$. That constant coordinate is not in $L_2[0,\infty)$; infinite-horizon stability theorems need the appropriate centering or affine/dissipativity treatment. Nonzero values $F(0)$ likewise require care when an $L_2\to L_2$ operator is required. All finite-horizon statements use locally square-integrable signals.

\section{Fixed saturation: QCs and IQCs}
Fix finite $\ell_i<u_i$, slopes $\kappa_i>0$, and offsets $\beta_i$. Define
\begin{equation}\label{eq:sat}
 \sigma_i(v_i)=\sat_{[\ell_i,u_i]}(\kappa_i v_i+\beta_i),\qquad
 \Sigma(v)=\col(\sigma_1(v_1),\ldots,\sigma_n(v_n)).
\end{equation}
The paper's repeated-function convention is recovered when all parameters are identical. The unit-slope, zero-offset saturation has $\kappa_i=1$, $\beta_i=0$. Parameters are fixed and known; for a parameter-uncertain family, intersect the valid multiplier cones over all admissible parameter choices.

\subsection{ReLU representation and a sufficient lifted multiplier family}
With $K=\diag(\kappa_i)$ and vector offsets and limits,
\begin{equation}\label{eq:satlift}
 a=\begin{bmatrix}Kv+\beta-\ell\\Kv+\beta-u\end{bmatrix},\qquad
 z=\Phi_{2n}(a),\qquad w=\ell+Jz,\qquad J=[I\ -I].
\end{equation}
This is an exact realization of $\Sigma$, not an approximation. Let $\xi=\col(v,z,1)$ and define the constant matrix $B$ by $B\xi=\col(a,z)$. Then
\begin{equation}\label{eq:lifted-sat-qc}
 \xi^TB^TQB\xi\ge0\qquad\forall Q\in\mc(\Phi_{2n}).
\end{equation}
In particular, one can use the full repeated-ReLU cone from Theorem 1 of~\cite{relu}, including its coupled constraints. Together with $w=\ell+Jz$, this is a lifted IQC description of saturation. Eliminating $z$ by a sufficient S-procedure may lose information; \eqref{eq:lifted-sat-qc} alone is not claimed to characterize every QC on $(v,w)$.

Simple additional affine QCs include
\begin{equation}\label{eq:amplitude}
 (w-\ell)^T\Lambda(u-w)\ge0,\quad\Lambda\in D^n,\ \Lambda\succeq0,
 \qquad
 \begin{bmatrix}w-\ell\\u-w\end{bmatrix}^TP
 \begin{bmatrix}w-\ell\\u-w\end{bmatrix}\ge0,\quad P\in\COP^{2n}.
\end{equation}
The first follows from the second for a suitable entrywise nonnegative $P$. Both are statements about the actual saturation outputs.

\subsection{A complete copositive characterization}
The following elementary homogenization replaces the orthant decomposition used for ReLU in~\cite{relu}.

\begin{lemma}[Polyhedral graph test]\label{lem:poly}
Let a nonempty pointed polyhedron $\mathcal P\subset\R^d$ have a finite representation
\[
 \mathcal P=\conv\{p_1,\ldots,p_r\}+\cone\{d_1,\ldots,d_s\}.
\]
Define its lifted generator
\begin{equation}\label{eq:generator}
 G(\mathcal P)=
 \begin{bmatrix}p_1&\cdots&p_r&d_1&\cdots&d_s\\1&\cdots&1&0&\cdots&0\end{bmatrix}.
\end{equation}
For $N\in\Ss^{d+1}$,
\[
 \col(p,1)^TN\col(p,1)\ge0\ \forall p\in\mathcal P
 \quad\Longleftrightarrow\quad
 G(\mathcal P)^TNG(\mathcal P)\in\COP^{r+s}.
\]
\end{lemma}
\begin{proof}
Every lifted point $\col(p,1)$ is a nonnegative combination of generator columns whose vertex weights sum to one. Conversely, a nonnegative combination with positive last coordinate is a positive scaling of such a lifted point. If the last coordinate is zero, add a vanishing positive multiple of any vertex column and use continuity of the quadratic form. This also proves the necessity of testing recession directions, not only vertices.
\end{proof}

Set
\[
 a_i^-=(\ell_i-\beta_i)/\kappa_i,\qquad
 a_i^+=(u_i-\beta_i)/\kappa_i.
\]
For $s\in\{-,0,+\}^n$, let $\mathcal I_s$ be the Cartesian product of
\[
 (-\infty,a_i^-],\quad[a_i^-,a_i^+],\quad[a_i^+,\infty)
\]
according to $s_i$. On $\mathcal I_s$, $\Sigma(v)=A_sv+b_s$, with rows respectively $\ell_i$, $\kappa_iv_i+\beta_i$, or $u_i$. Define
\[
 \mathcal P_s=\{\col(v,w):v\in\mathcal I_s,\ w=A_sv+b_s\},\qquad G_s=G(\mathcal P_s).
\]
These generators are explicit: choose all input vertices by selecting either endpoint in each middle interval and the finite endpoint in each tail interval; lift each as $\col(v,A_sv+b_s,1)$. Add $\col(-e_i,0,0)$ for each lower tail and $\col(e_i,0,0)$ for each upper tail. If $m_s$ coordinates are middle intervals, there are $r_s=2^{m_s}+n-m_s$ columns.

\begin{theorem}[Complete static saturation QCs]\label{thm:satqc}
Define
\begin{align}
 \widetilde{\mc}(\Sigma)
 &=\{N\in\Ss^{2n+1}:G_s^TNG_s\in\COP^{r_s}\ \forall s\in\{-,0,+\}^n\},\label{eq:sat-aff-cone}\\
 \mc(\Sigma)
 &=\{M\in\Ss^{2n}:E^TME\in\widetilde{\mc}(\Sigma)\}.\label{eq:sat-cone}
\end{align}
Then $N\in\widetilde{\mc}(\Sigma)$ if and only if \eqref{eq:affqc} holds on the saturation graph. Also $\Sigma\in QC(M)$ if and only if $M\in\mc(\Sigma)$. Every such multiplier gives the corresponding hard IQC.
\end{theorem}
\begin{proof}
The saturation graph is exactly the union of the $3^n$ polyhedra $\mathcal P_s$. Apply Lemma~\ref{lem:poly} to each graph piece. Substitute $N=E^TME$ for the homogeneous statement and integrate for the hard IQC.
\end{proof}

\section{Fixed saturation: incremental QCs and IIQCs}
Set $\delta v=\bar v-\hat v$ and $\delta w=\Sigma(\bar v)-\Sigma(\hat v)$. Then
\begin{equation}\label{eq:sat-slope}
 K^{-1}\delta w=D\delta v,\qquad D\in D^n,\quad0\preceq D\preceq I.
\end{equation}
Let $T_K=\diag(I,K^{-1})$.

\begin{theorem}[Complete homogeneous incremental saturation QCs]\label{thm:satinc}
The complete cone is
\begin{equation}\label{eq:sat-inc-complete}
 \boxed{\mc^{inc}(\Sigma)=
 \{T_K^TQT_K:Q\in\mc^{inc}(\Phi_n)\}.}
\end{equation}
Consequently every $M$ in this cone supplies a hard static IIQC
\[
 \int_0^T\col(\delta v,\delta w)^TM\col(\delta v,\delta w)\,dt\ge0.
\]
For unit slopes, the complete homogeneous incremental multiplier cone is exactly the paper's incremental repeated-ReLU cone, independently of the fixed finite saturation limits and offsets.
\end{theorem}
\begin{proof}
First consider unit slopes and zero offsets. The set of repeated-ReLU increments is precisely
\[
 \mathcal C=\{\col(v,w):w=Dv,\ D\in D^n,\ 0\preceq D\preceq I\}.
\]
The inclusion into $\mathcal C$ follows from the ReLU slope restriction. To prove the reverse inclusion and the saturation claim simultaneously, take any $(v,w)\in\mathcal C$ and choose $\tau>0$ with $\tau|v_i|<u_i-\ell_i$ for every nonzero $v_i$. For $v_i\ge0$, choose
\[
 \bar v_i=\ell_i+\tau w_i,\qquad
 \hat v_i=\ell_i+\tau(w_i-v_i).
\]
For $v_i<0$, choose
\[
 \bar v_i=\ell_i+\tau(v_i-w_i),\qquad
 \hat v_i=\ell_i-\tau w_i.
\]
One argument is below the lower kink and the other is in the linear segment. Their input/output differences equal $(\tau v_i,\tau w_i)$. The same construction at the ReLU kink, with no upper clipping, realizes $(v,w)$ exactly for ReLU. Every saturation increment lies in $\mathcal C$, and every direction in $\mathcal C$ occurs as a scaled saturation increment. Homogeneity of the quadratic form proves equality of the multiplier cones. For general $K,\beta$, use the input thresholds $a_i^-$ and normalized output differences $K^{-1}\delta w$; the linear input widths are $(u_i-\ell_i)/\kappa_i$. This yields the congruence in \eqref{eq:sat-inc-complete}.
\end{proof}

The paper's diagonal sector cone is one available subset. A coupled member for two unit-slope channels is
\begin{equation}\label{eq:cross}
 (\delta v_1-\delta v_2)^2+2\delta w_1\delta w_2\ge0,
 \qquad
 M_\times=\begin{bmatrix}1&-1&0&0\\-1&1&0&0\\0&0&0&1\\0&0&1&0\end{bmatrix}.
\end{equation}
It is valid for arbitrary componentwise slope-$[0,1]$ functions. The paper derives this in Proposition 2 and builds a coupled multiplier family in Theorem 4. Conic combinations of these and the existing hard dynamic sigmoid IQCs remain hard. This does not establish dominance over the existing dynamic family or feasibility for a particular closed loop.

\begin{remark}
Exact saturation levels can be exploited by the affine graph QCs or by a two-trajectory lift retaining absolute inputs. Equation~\eqref{eq:sat-inc-complete} shows why adding only homogeneous static quadratic forms of $(\delta v,\delta w)$ cannot retain the absolute saturation levels. A complete cone of valid QCs need not uniquely identify a function: the sign-invariance issue discussed in~\cite{relu} still applies to homogeneous graph QCs.
\end{remark}

\section{The class between two saturation envelopes: QCs and IQCs}
Fix two maps $\underline\Sigma,\overline\Sigma$ of the form \eqref{eq:sat}, possibly with different parameters, and assume $\underline\Sigma(v)\le\overline\Sigma(v)$ componentwise for every $v$. Define
\begin{equation}\label{eq:env-class}
 \Fenv=\{F:\ F(v)=\col(f_i(v_i)),\quad
 \underline\sigma_i(r)\le f_i(r)\le\overline\sigma_i(r)\quad\forall r,i\}.
\end{equation}
This definition imposes no monotonicity or slope restriction. Requiring a common repeated scalar function is an optional further specialization when the envelopes are identical across channels. The results below for the unrestricted class also hold when the selections are required to be continuous.

For $w=F(v)$, write
\begin{equation}\label{eq:env-slack}
 s=\begin{bmatrix}w-\underline\Sigma(v)\\\overline\Sigma(v)-w\end{bmatrix}\ge0.
\end{equation}
Thus for every $P\in\COP^{2n}$,
\begin{equation}\label{eq:env-cop}
 s^TPs\ge0.
\end{equation}
A useful special case is
\begin{equation}\label{eq:env-sector}
 (w-\underline\Sigma(v))^T\Lambda(\overline\Sigma(v)-w)\ge0,
 \qquad\Lambda\in D^n,\quad\Lambda\succeq0.
\end{equation}

\begin{proposition}[Exact membership in the sandwich class]\label{prop:membership}
With the envelopes fixed and ordered as above, a componentwise function $F$ belongs to $\Fenv$ if and only if
\begin{equation}\label{eq:membership-iff}
 (f_i(v_i)-\underline\sigma_i(v_i))
 (\overline\sigma_i(v_i)-f_i(v_i))\ge0\quad\forall v_i,i.
\end{equation}
Equivalently, \eqref{eq:env-sector} holds for every $v$ and every nonnegative diagonal $\Lambda$. Equivalently, its hard integral inequality holds for every input trajectory, every $T\ge0$, and every such $\Lambda$.
\end{proposition}
\begin{proof}
For ordered real numbers $l\le u$, $(y-l)(u-y)\ge0$ holds exactly when $l\le y\le u$. This includes $l=u$. Select each diagonal coordinate separately for the vector equivalence. Pointwise validity implies hard integral validity. For the converse, use a constant input and any positive finite horizon.
\end{proof}
This is an exact \emph{lifted} membership description once the envelope realizations are retained. It does not assert that a sufficient S-procedure eliminating those realizations is exact, nor that completeness of the valid homogeneous multiplier cone uniquely identifies the original function class.

\subsection{ReLU lift using the paper's multiplier cone}
Stack the two envelope realizations, so
\begin{equation}\label{eq:env-lift}
 a=Av+b\in\R^{4n},\qquad z=\Phi_{4n}(a),\qquad
 \underline\Sigma(v)=\ell_-+C_-z,\quad
 \overline\Sigma(v)=\ell_++C_+z.
\end{equation}
Explicitly, the four blocks of $a$ are $K_-v+\beta_--\ell_-$, $K_-v+\beta_--u_-$, $K_+v+\beta_+-\ell_+$, and $K_+v+\beta_+-u_+$; $C_-=[I,-I,0,0]$, $C_+=[0,0,I,-I]$.
Set $\xi=\col(v,w,z,1)$ and define constant maps
\[
 B_\Phi\xi=\col(Av+b,z),\qquad
 B_s\xi=\col(w-\ell_--C_-z,\ell_++C_+z-w).
\]
For every $F\in\Fenv$, the sufficient lifted QC family is
\begin{equation}\label{eq:env-qc-family}
 \boxed{\xi^T\left(B_\Phi^TQB_\Phi+B_s^TPB_s\right)\xi\ge0,
 \quad Q\in\mc(\Phi_{4n}),\quad P\in\COP^{2n}.}
\end{equation}
It retains both the repeated-ReLU structure of the envelopes and the two-sided slack constraints. It is a quadratic constraint on a lifted graph, with auxiliary nonlinear outputs $z$; it is not yet a multiplier solely on $(v,w)$. Integration gives its hard lifted IQC. In a plant interconnection, add the auxiliary outputs explicitly or perform a justified elimination.

\subsection{Complete QCs on the original graph}
Partition each scalar input axis at all breakpoints of its two envelopes. The resulting Cartesian cells $\mathcal I_j$ make both envelopes affine:
\[
 \underline\Sigma(v)=A_j^-v+b_j^-,\qquad
 \overline\Sigma(v)=A_j^+v+b_j^+\quad(v\in\mathcal I_j).
\]
Define the polyhedral strips
\begin{equation}\label{eq:strip}
 \mathcal P_j=\{\col(v,w):v\in\mathcal I_j,\quad
 A_j^-v+b_j^-\le w\le A_j^+v+b_j^+\},\qquad G_j=G(\mathcal P_j).
\end{equation}
Empty pieces can be discarded. In one scalar cell with finite endpoints, the strip vertices are its lower and upper envelope values at each endpoint (remove duplicates). In a tail cell both envelopes are constant: use the two values at its finite endpoint and the input recession ray $(\pm1,0)$. Cartesian products give the vector vertices; recession rays are embedded one coordinate at a time.

\begin{theorem}[Complete QCs for the envelope class]\label{thm:envqc}
With $r_j$ the number of columns of $G_j$, define
\begin{align}
 \widetilde{\mc}(\Fenv)
 &=\{N\in\Ss^{2n+1}:G_j^TNG_j\in\COP^{r_j}\ \forall j\},\label{eq:env-aff-complete}\\
 \mc(\Fenv)
 &=\{M\in\Ss^{2n}:E^TME\in\widetilde{\mc}(\Fenv)\}.
\end{align}
These are exactly the affine and homogeneous graph multipliers valid for every $F\in\Fenv$. They give hard IQCs by integration.
\end{theorem}
\begin{proof}
The union of the strips is precisely the union of the graphs of all componentwise selections in \eqref{eq:env-class}. Every allowed value at one input can be attained by a selection, including a continuous selection: choose a constant mixing coefficient between the continuous lower and upper envelopes in each channel. Apply Lemma~\ref{lem:poly} on each piece.
\end{proof}

\section{The envelope class: incremental QCs and IIQCs}
Pointwise envelope ordering does not imply
\[
 \underline\Sigma(\bar v)-\underline\Sigma(\hat v)
 \le F(\bar v)-F(\hat v)
 \le\overline\Sigma(\bar v)-\overline\Sigma(\hat v).
\]
The valid bounds are
\begin{equation}\label{eq:two-env}
 \underline\Sigma(\bar v)-\overline\Sigma(\hat v)
 \le\delta w\le
 \overline\Sigma(\bar v)-\underline\Sigma(\hat v).
\end{equation}
In particular, with $\underline b_\Delta$ and $\overline b_\Delta$ the left and right bounds,
\begin{equation}\label{eq:two-env-qc}
 (\delta w-\underline b_\Delta)^T\Lambda(\overline b_\Delta-\delta w)\ge0,
 \qquad\Lambda\in D^n,\quad\Lambda\succeq0.
\end{equation}
This is a lifted two-trajectory QC depending on absolute inputs, not generally a homogeneous QC solely in $(\delta v,\delta w)$.

\subsection{Complete copositive incremental test}
For each pair of strip pieces, form the polyhedron
\[
 \mathcal P_{jk}^{(2)}=
 \{\col(\bar v,\hat v,\bar w,\hat w):
 \col(\bar v,\bar w)\in\mathcal P_j,\quad
 \col(\hat v,\hat w)\in\mathcal P_k\}.
\]
Its vertices are all pairs of vertices of $\mathcal P_j,\mathcal P_k$; its recession rays vary one factor at a time. Let $G_{jk}^{(2)}$ be its lifted generator. Use exactly the difference map from the proof of Theorem 3 in~\cite{relu}:
\begin{equation}\label{eq:R}
 R=\begin{bmatrix}I&-I&0&0\\0&0&I&-I\end{bmatrix},\qquad
 H_{jk}=[R\ \ 0]G_{jk}^{(2)}.
\end{equation}
\begin{theorem}[Complete homogeneous incremental envelope QCs]\label{thm:envinc}
With $r_{jk}$ the number of columns of $G_{jk}^{(2)}$,
\begin{equation}\label{eq:env-inc-complete}
 \boxed{\mc^{inc}(\Fenv)=
 \{M\in\Ss^{2n}:H_{jk}^TMH_{jk}\in\COP^{r_{jk}}\quad\forall j,k\}.}
\end{equation}
Every $M$ in this cone gives a hard IIQC on the differences.
\end{theorem}
\begin{proof}
Testing all product-strip points is sufficient for every same-function trajectory pair. It is also necessary for the unrestricted selection class. If the two scalar inputs in a channel are distinct, any two allowed outputs can be attained by a single selection: interpolate their mixing coefficients between the continuous envelopes. Such a selection can be continuous. If the two inputs coincide but the chosen outputs differ, the product strips contain a point that is not a same-function pair. Approximate it with distinct inputs and feasible outputs converging to the prescribed values; continuity of the envelopes makes this possible. Validity of the quadratic form extends to that limit. Thus the product strips are the closure of the union of same-function two-point graphs. Apply Lemma~\ref{lem:poly} with the pullback $[R\ 0]^TM[R\ 0]$. For a common repeated scalar selection, perturb all finitely many scalar inputs to be distinct and use the same argument.
\end{proof}

\begin{proposition}[A separated sandwich alone has trivial difference QCs]\label{prop:collapse}
Suppose each channel has some $r_i$ with $\underline\sigma_i(r_i)<\overline\sigma_i(r_i)$. Then
\begin{equation}\label{eq:trivial}
 \boxed{\mc^{inc}(\Fenv)=\Ss_+^{2n}.}
\end{equation}
This also holds for the class of all continuous selections. If the envelopes coincide identically, this conclusion need not hold; the class then reduces to the prescribed function.
\end{proposition}
\begin{proof}
Choose in each channel a point strictly inside the envelope strip. For any direction $(a,b)$ with all $a_i\ne0$, two sufficiently close inputs and outputs with differences $(\tau a,\tau b)$ lie in these strict strips. A single continuous selection can attain the two values in each channel. If a homogeneous incremental QC is valid for every selection, divide its inequality by $\tau^2$ to obtain $\col(a,b)^TM\col(a,b)\ge0$. Directions with some $a_i=0$ follow by density. This proves $M\succeq0$; the converse is immediate. For a common scalar selection and identical envelopes, choose distinct interior basepoints and use a continuous interpolation through the finitely many assigned values.
\end{proof}

\subsection{The useful sigmoid subclass: sandwich and slope together}
For normalized sigmoid slopes define
\begin{equation}\label{eq:sigclass}
 \Fenv^{[0,1]}=\{F\in\Fenv:0\le[f_i(r)-f_i(s)]/(r-s)\le1
 \ \forall r\ne s,i\}.
\end{equation}
This is strictly more information than the sandwich alone. Every $F$ in this subclass satisfies all $Q_\Delta\in\mc^{inc}(\Phi_n)$, not only the diagonal sector family, as follows from the attainable increment set $\mathcal C$ in the proof of Theorem~\ref{thm:satinc}. It also satisfies the one-copy envelope QCs and the two-copy lifted constraints below.

Take $\eta=\col(\bar v,\hat v,\bar w,\hat w,\bar z,\hat z,1)$, with the envelope realization \eqref{eq:env-lift} in each copy. Define constant linear maps by
\begin{align*}
 L_\Delta\eta&=\col(\bar v-\hat v,\bar w-\hat w),\\
 L_\Phi\eta&=\col(A\bar v+b,A\hat v+b,\bar z,\hat z),\\
 L_s\eta&=\col(\bar w-\ell_--C_-\bar z,\ell_++C_+\bar z-\bar w,
                 \hat w-\ell_--C_-\hat z,\ell_++C_+\hat z-\hat w).
\end{align*}
The following is a valid sufficient lifted family:
\begin{equation}\label{eq:subclass-lift}
 \boxed{\eta^TQ_{\rm lift}\eta\ge0,\qquad
 Q_{\rm lift}=L_\Phi^TQL_\Phi+L_s^TPL_s+L_\Delta^TQ_\Delta L_\Delta,}
\end{equation}
where $Q\in\mc(\Phi_{8n})$, $P\in\COP^{4n}$, and $Q_\Delta\in\mc^{inc}(\Phi_n)$. The two trajectories are stacked as one repeated-ReLU graph, so cross-copy ReLU QCs are allowed. Nonnegative conic combinations and hard integration preserve validity. One may also add the explicit constraint \eqref{eq:two-env-qc} or incremental ReLU QCs for the auxiliary envelope channels.

To obtain a sufficient standard multiplier solely on the original differences, find these multipliers and $M\in\Ss^{2n}$ such that
\begin{equation}\label{eq:eliminate}
 \boxed{L_\Delta^TML_\Delta-Q_{\rm lift}\succeq0.}
\end{equation}
Then $F\in QC^{inc}(M)$ for every $F\in\Fenv^{[0,1]}$, hence $M$ gives a hard IIQC. Equation~\eqref{eq:eliminate} is sufficient and may be conservative; no completeness or strict improvement is asserted for this elimination. Alternatively, retain the full two-copy lift in a suitable dissipation analysis rather than identifying it with a difference-only IIQC.

\section{Concrete global sigmoid sandwiches}
\subsection{Horizontally shifted ramps with small vertical margins}
The geometry of the proposed sketch is obtained by putting the upper envelope slightly above both sigmoid asymptotes and the lower envelope slightly below them. Both ramps can have the same slope and height, while being displaced horizontally. For
\[
 S(q)=\frac M2\left(1+\tanh\frac{2(q-c)}M\right),\qquad M>0,
\]
choose a margin $0<\varepsilon<M/2$ and a ramp slope $0<k\le1$. Define
\begin{align*}
 q_\varepsilon&=c+\frac M4\log\frac{\varepsilon}{M-\varepsilon},
 &h&=\frac M2\operatorname{arcosh}\frac1{\sqrt k},\\
 S_R&=\frac M2(1+\sqrt{1-k}),
 &a_*&=\min\left\{q_\varepsilon,\ c+h-\frac{S_R-\varepsilon}{k}\right\},
 &b_*&=2c-a_*-\frac Mk.
\end{align*}
\begin{proposition}[Necessary and sufficient horizontal shifts]\label{prop:shift-iff}
For the two envelopes
\begin{equation}\label{eq:shifted-envelope}
 \underline\sigma_b(q)=-\varepsilon+\sat_{[0,M]}(k(q-b)),\qquad
 \overline\sigma_a(q)=\varepsilon+\sat_{[0,M]}(k(q-a)),
\end{equation}
the exact conditions are
\begin{align}
 S(q)\le\overline\sigma_a(q)\quad\forall q
 &\quad\Longleftrightarrow\quad a\le a_*,\label{eq:upper-iff}\\
 \underline\sigma_b(q)\le S(q)\quad\forall q
 &\quad\Longleftrightarrow\quad b\ge b_* .\label{eq:lower-iff}
\end{align}
Thus $\underline\sigma_b\le S\le\overline\sigma_a$ globally if and only if $a\le a_*$ and $b\ge b_*$. The choices $a=a_*$, $b=b_*$ give the tightest horizontal positions within this fixed $(\varepsilon,k)$ family.
\end{proposition}

\begin{proof}
For $q\le q_\varepsilon$, $S(q)\le\varepsilon\le\overline\sigma_a(q)$. For $q>q_\varepsilon$, the upper bound is equivalent to
\[
 a\le q-\frac{S(q)-\varepsilon}{k}.
\]
Consequently the exact upper threshold is $\inf_{q>q_\varepsilon}[q-(S(q)-\varepsilon)/k]$. The right side tends to infinity as $q\to\infty$. Its derivative is $1-S'(q)/k$; for $k<1$ its only possible interior minimum is at $q=c+h$, where $S(q)=S_R$. Its boundary limit is $q_\varepsilon$, giving precisely $a_*$. For $k=1$ the function is nondecreasing, and the same formula gives the boundary infimum. For the lower bound, use $S(q)=M-S(2c-q)$ and
\[
 M-\underline\sigma_b(2c-q)
 =\varepsilon+\sat_{[0,M]}\bigl(k(q-(2c-b-M/k))\bigr).
\]
The upper-bound equivalence gives $2c-b-M/k\le a_*$, which is exactly $b\ge b_*$. This proves necessity as well as sufficiency.
\end{proof}

The envelopes have equal ramp slope $k$ and equal vertical span $M$; their plateaus are $(-\varepsilon,M-\varepsilon)$ and $(\varepsilon,M+\varepsilon)$. Thus they look nearly equal at the top and bottom when $\varepsilon$ is small, but the strict vertical margins are retained in the exact statement. Smaller margins need not produce a smaller central gap: the horizontal shifts can grow.

For example, $M=1$, $c=0$, $\varepsilon=0.01$, and $k=0.5$ give $a_*\simeq-1.2464199877$ and $b_*\simeq-0.7535800123$. Take the threshold choices and write $\underline\sigma=\underline\sigma_{b_*}$, $\overline\sigma=\overline\sigma_{a_*}$. The red upper ramp then starts earlier than the blue lower ramp, and their top levels are $1.01$ and $0.99$. If an exactly zero lower floor is desired, replace the lower envelope by
\begin{equation}\label{eq:zero-floor}
 \underline\sigma_0(q)=\max\{0,\underline\sigma(q)\}
 =\sat_{[0,M-\varepsilon]}(k(q-b_*)-\varepsilon)\le S(q).
\end{equation}
This retains a two-ReLU realization. Since $S\ge0$, the equivalence \eqref{eq:lower-iff} also holds if $\underline\sigma_b$ is replaced by $\max\{0,\underline\sigma_b\}$. With or without this clipping, the envelope QCs and IIQC constructions above apply unchanged. In particular, the vertical margins do not establish a slope restriction for every function inside the sandwich; the actual sigmoid's known slope bound is still additional information.

\subsection{An alternative construction from a uniform ramp error}
For
\[
 S(q)=\frac{M}{2}\left(1+\tanh\frac{2(q-c)}{M}\right),\qquad M>0,
\]
define
\[
 p(q)=\sat_{[0,M]}(q-c+M/2),\qquad
 \varepsilon=\frac M2(1-\tanh1).
\]
On the ramp, $(S-p)'=S'-1\le0$. On the two tails, $S$ is increasing and $p$ is constant. Symmetry about $(c,M/2)$ therefore gives $|S-p|\le\varepsilon$, with equality at the ramp endpoints. Hence
\begin{equation}\label{eq:concrete}
 \underline\sigma(q)=\sat_{[0,M-\varepsilon]}(q-c+M/2-\varepsilon)
 \le S(q)\le
 \sat_{[\varepsilon,M]}(q-c+M/2+\varepsilon)=\overline\sigma(q).
\end{equation}
Also $0<S'(q)\le1$, so this sigmoid belongs to the subclass \eqref{eq:sigclass}. The two envelopes have different plateau levels; two finite-breakpoint saturations both using plateaus $0,M$ cannot provide a global lower and upper bound for a sigmoid strictly between $0,M$. Tighter envelopes or several piecewise linear segments can be substituted; the polyhedral construction and the ReLU-lift argument remain applicable.

\section{Computation and scope}
All complete characterizations here are exact copositive conditions. As in~\cite{relu}, a tractable sufficient SDP family replaces a required copositive matrix $C$ by
\[
 C=S+N,\qquad S\succeq0,\quad N=N^T\ge0\ \text{entrywise}.
\]
This replacement is an inner approximation of the admissible multiplier cone and can be conservative. The same relaxation can be applied to $P\in\COP$ in the lifted families. Equality constraints and all constant offsets must be retained. Multiplier validity, feasibility of a complementary plant/controller inequality, and interconnection well-posedness are distinct requirements. These derivations do not certify the workspace's fixed delayed closed loop.

The accompanying Python script checks the explicit saturation and envelope generators, their two-copy difference maps, the scaled-increment construction, the analytic sigmoid envelope, and a continuous sandwich selection violating a unit slope bound. These numerical checks verify the formulas and constructions, not universal copositivity or a closed-loop stability result.

\begin{thebibliography}{1}
\bibitem{relu} S. Vahedi Noori, B. Hu, G. Dullerud, and P. Seiler,
``A Complete Set of Quadratic Constraints for Repeated ReLU and Generalizations,''
arXiv:2407.06888v2, 22 August 2024.
\url{https://arxiv.org/abs/2407.06888v2}.
\end{thebibliography}
\end{document}
